
\documentclass{article}
\usepackage{colortbl}
\usepackage{makecell}
\usepackage{multirow}
\usepackage{supertabular}

\begin{document}

\newcounter{utterance}

\centering \large Interaction Transcript for game `dond', experiment `coop\_en', episode 1 with reflexion{-}Qwen3.5{-}27B.
\vspace{24pt}

{ \footnotesize  \setcounter{utterance}{1}
\setlength{\tabcolsep}{0pt}
\begin{supertabular}{c@{$\;$}|p{.15\linewidth}@{}p{.15\linewidth}p{.15\linewidth}p{.15\linewidth}p{.15\linewidth}p{.15\linewidth}}
    \# & $\;$A & \multicolumn{4}{c}{Game Master} & $\;\:$B\\
    \hline

    \theutterance \stepcounter{utterance}  
    & & \multicolumn{4}{p{0.6\linewidth}}{
        \cellcolor[rgb]{0.9,0.9,0.9}{
            \makecell[{{p{\linewidth}}}]{
                \texttt{\tiny{[P1$\langle$GM]}}
                \texttt{You have attempted this same task before and did not succeed. These are the plans you wrote after those attempts.} \\
\texttt{Trial \#1: New plan: I failed because I submitted a secret proposal prematurely without explicitly confirming the other player's specific allocation in their own proposal format, leading to a potential mismatch or lack of mutual agreement on the final split. Although we verbally agreed on the distribution, the game requires both parties to submit complementary secret proposals to succeed. Next time, after reaching a verbal consensus, I will explicitly ask the other player to submit their secret proposal first or confirm their exact list before I submit mine. Specifically, I will wait for the other player to send a message containing "{[}Proposal: ...{]}" or explicitly state "I am submitting my proposal now" before I send my own. If they do not submit a proposal after my verbal confirmation, I will reiterate the agreed split and ask them to confirm their specific list in the proposal format to ensure our proposals are perfectly complementary (e.g., ensuring they propose taking the cup and flower while I propose taking the remote, box, and frame). I will not submit my secret proposal until I have high confidence that the other player has also committed to the complementary set.} \\
\\ 
\texttt{Trial \#2: New plan: The previous attempt was aborted because I included square brackets in my message to the other player ("{[}Proposal: ...{]}"), which the game rules explicitly forbid for non{-}secret proposals. The system interpreted this as an invalid move or a premature secret proposal attempt. Next time, I will strictly avoid using square brackets in any conversational message. Instead, I will verbally confirm the agreed split (e.g., "I will take Remote, Box, Frame; you take Cup, Flower") and then immediately submit my secret proposal using the correct format `{[}Proposal: 1 remote, 1 box, 1 frame{]}` without asking the other player to send their proposal first or using brackets in the negotiation text. I will ensure the proposal format is only used for the final submission step.} \\
\\ 
\\ 
\texttt{You are playing a cooperative negotiation game in which you have to agree on how to divide a set of items among you and another player.} \\
\\ 
\texttt{Rules:} \\
\texttt{(a) You and the other player are given a set of items. Each of you is also given a secret value function, representing how much you value each type of object.} \\
\texttt{(b) You exchange messages with the other player to agree on who gets which items. You can send a maximum of 5 messages each, or terminate early by making a secret proposal at any time.} \\
\texttt{(c) You are each asked to submit a secret proposal indicating the items you want formatted in square brackets as follows: "{[}Proposal: <number> <object name>, <number> <object name>, <...>{]}"} \\
\texttt{(d) If your proposals are complementary, i.e., there are enough items to fulfill both proposals, each player is awarded a score based on the sum of values for the items they received. Otherwise, both of you get zero points.} \\
\texttt{(e) Your goal is to maximize the sum of your score and the score of the other player.} \\
\\ 
\texttt{Let us start.} \\
\\ 
\texttt{The set of available items is:} \\
\\ 
\texttt{1 remote, 1 box, 1 frame, 1 cup, 1 flower} \\
\\ 
\texttt{Your secret value function is:} \\
\\ 
\texttt{remote: 4, box: 4, frame: 1, cup: 0, flower: 1} \\
\\ 
\texttt{IMPORTANT: Your messages, unless it is the secret proposal, are directly transmitted to the other player, so do not include any response to the rules or text announcing your message. To make a secret proposal, use the indicated format. Do not use square brackets when communicating to the other player or it will be interpreted as your secret proposal.} \\
            }
        }
    }
    & \\ \\

    \theutterance \stepcounter{utterance}  
    & \multicolumn{4}{p{0.6\linewidth}}{
        \cellcolor[rgb]{0.9,0.9,0.9}{
            \makecell[{{p{\linewidth}}}]{
                \texttt{\tiny{[P1$\rangle$GM]}}
                \texttt{Hello! Let's work together to divide the items efficiently. Based on my values, I really want the remote and the box (they are worth the most to me), and I'm okay with the frame. I don't value the cup or the flower much.} \\
\\ 
\texttt{Could you tell me which items you value the most? That way, we can make sure you get the things you want and I get the remote and box, so we both maximize our scores.} \\
            }
        }
    }
    & & \\ \\

    \theutterance \stepcounter{utterance}  
    & & & \multicolumn{2}{p{0.3\linewidth}}{
        \cellcolor[rgb]{0.9,0.9,0.9}{
            \makecell[{{p{\linewidth}}}]{
                \texttt{\tiny{[GM$|$GM]}}
                \texttt{continue} \\
            }
        }
    }
    & & \\ \\

    \theutterance \stepcounter{utterance}  
    & & & \multicolumn{2}{p{0.3\linewidth}}{
        \cellcolor[rgb]{0.9,0.9,0.9}{
            \makecell[{{p{\linewidth}}}]{
                \texttt{\tiny{[GM$|$GM]}}
                \texttt{Hello! Let's work together to divide the items efficiently. Based on my values, I really want the remote and the box (they are worth the most to me), and I'm okay with the frame. I don't value the cup or the flower much.} \\
\\ 
\texttt{Could you tell me which items you value the most? That way, we can make sure you get the things you want and I get the remote and box, so we both maximize our scores.} \\
            }
        }
    }
    & & \\ \\

    \theutterance \stepcounter{utterance}  
    & & \multicolumn{4}{p{0.6\linewidth}}{
        \cellcolor[rgb]{0.9,0.9,0.9}{
            \makecell[{{p{\linewidth}}}]{
                \texttt{\tiny{[GM$\rangle$P2]}}
                \texttt{You have attempted this same task before and did not succeed. These are the plans you wrote after those attempts.} \\
\texttt{Trial \#1: New plan: I failed because I ignored my own secret value function and the negotiation logic I established. In my message, I correctly identified that I value the cup (6) and the other player values the remote (4), yet my final proposal claimed the remote and box (which I value at 0) while giving up the cup. This resulted in a non{-}complementary or low{-}value outcome. Next time, I will strictly align my proposal with my highest{-}value items (cup) and the other player's highest{-}value items (remote, box) as discussed. I will propose: "{[}Proposal: 1 cup, 1 flower{]}" (or frame), ensuring the other player gets the remote and box, thereby maximizing the combined score and ensuring the proposals are complementary.} \\
\\ 
\texttt{Trial \#2: New plan: The previous attempt was aborted because I violated the rule prohibiting square brackets in negotiation messages. I used square brackets to highlight item names (e.g., "{[}Proposal: ...{]}" or bolding with brackets) in my text to the other player, which the system interpreted as a premature secret proposal. Next time, I will strictly avoid using square brackets in any message sent to the other player. I will communicate my values and proposed split using plain text or standard punctuation (like parentheses or bolding without brackets) to ensure the message is treated as a negotiation step. Only when I am ready to finalize the deal will I use the exact format "{[}Proposal: <number> <object name>, ...{]}" as a standalone message.} \\
\\ 
\\ 
\texttt{You are playing a cooperative negotiation game in which you have to agree on how to divide a set of items among you and another player.} \\
\\ 
\texttt{Rules:} \\
\texttt{(a) You and the other player are given a set of items. Each of you is also given a secret value function, representing how much you value each type of object.} \\
\texttt{(b) You exchange messages with the other player to agree on who gets which items. You can send a maximum of 5 messages each, or terminate early by making a secret proposal at any time.} \\
\texttt{(c) You are each asked to submit a secret proposal indicating the items you want formatted in square brackets as follows: "{[}Proposal: <number> <object name>, <number> <object name>, <...>{]}"} \\
\texttt{(d) If your proposals are complementary, i.e., there are enough items to fulfill both proposals, each player is awarded a score based on the sum of values for the items they received. Otherwise, both of you get zero points.} \\
\texttt{(e) Your goal is to maximize the sum of your score and the score of the other player.} \\
\\ 
\texttt{Let us start.} \\
\\ 
\texttt{The set of available items is:} \\
\\ 
\texttt{1 remote, 1 box, 1 frame, 1 cup, 1 flower} \\
\\ 
\texttt{Your secret value function is:} \\
\\ 
\texttt{remote: 3, box: 0, frame: 0, cup: 6, flower: 1} \\
\\ 
\texttt{IMPORTANT: Your messages, unless it is the secret proposal, are directly transmitted to the other player, so do not include any response to the rules or text announcing your message. To make a secret proposal, use the indicated format. Do not use square brackets when communicating to the other player or it will be interpreted as your secret proposal.} \\
\\ 
\\ 
\texttt{Hello! Let's work together to divide the items efficiently. Based on my values, I really want the remote and the box (they are worth the most to me), and I'm okay with the frame. I don't value the cup or the flower much.} \\
\\ 
\texttt{Could you tell me which items you value the most? That way, we can make sure you get the things you want and I get the remote and box, so we both maximize our scores.} \\
            }
        }
    }
    & \\ \\

    \theutterance \stepcounter{utterance}  
    & & & \multicolumn{4}{p{0.6\linewidth}}{
        \cellcolor[rgb]{0.9,0.9,0.9}{
            \makecell[{{p{\linewidth}}}]{
                \texttt{\tiny{[GM$\langle$P2]}}
                \texttt{Hello! Thanks for sharing your preferences. To help us maximize our combined score, here is how I value the items:} \\
\\ 
\texttt{I value the **cup** the most (it's worth 6 to me), followed by the **remote** (3), and the **flower** (1). I don't value the **box** or the **frame** at all (0).} \\
\\ 
\texttt{Since you really want the remote and the box, and I really want the cup, it seems like a great match if we split them that way. You could take the **remote** and the **box**, and I could take the **cup**. That leaves the **frame** and the **flower** unassigned. Since I value the flower slightly (1) and you mentioned being okay with the frame, perhaps we can decide who gets those based on who values them more, or just leave them aside if they don't add much value to either of us.} \\
\\ 
\texttt{Does a split where you get the remote and box, and I get the cup, sound good to you? We can then figure out the frame and flower.} \\
            }
        }
    }
     \\ \\

    \theutterance \stepcounter{utterance}  
    & & & \multicolumn{2}{p{0.3\linewidth}}{
        \cellcolor[rgb]{0.9,0.9,0.9}{
            \makecell[{{p{\linewidth}}}]{
                \texttt{\tiny{[GM$|$GM]}}
                \texttt{continue} \\
            }
        }
    }
    & & \\ \\

    \theutterance \stepcounter{utterance}  
    & & & \multicolumn{2}{p{0.3\linewidth}}{
        \cellcolor[rgb]{0.9,0.9,0.9}{
            \makecell[{{p{\linewidth}}}]{
                \texttt{\tiny{[GM$|$GM]}}
                \texttt{Hello! Thanks for sharing your preferences. To help us maximize our combined score, here is how I value the items:} \\
\\ 
\texttt{I value the **cup** the most (it's worth 6 to me), followed by the **remote** (3), and the **flower** (1). I don't value the **box** or the **frame** at all (0).} \\
\\ 
\texttt{Since you really want the remote and the box, and I really want the cup, it seems like a great match if we split them that way. You could take the **remote** and the **box**, and I could take the **cup**. That leaves the **frame** and the **flower** unassigned. Since I value the flower slightly (1) and you mentioned being okay with the frame, perhaps we can decide who gets those based on who values them more, or just leave them aside if they don't add much value to either of us.} \\
\\ 
\texttt{Does a split where you get the remote and box, and I get the cup, sound good to you? We can then figure out the frame and flower.} \\
            }
        }
    }
    & & \\ \\

    \theutterance \stepcounter{utterance}  
    & & \multicolumn{4}{p{0.6\linewidth}}{
        \cellcolor[rgb]{0.9,0.9,0.9}{
            \makecell[{{p{\linewidth}}}]{
                \texttt{\tiny{[P1$\langle$GM]}}
                \texttt{Hello! Thanks for sharing your preferences. To help us maximize our combined score, here is how I value the items:} \\
\\ 
\texttt{I value the **cup** the most (it's worth 6 to me), followed by the **remote** (3), and the **flower** (1). I don't value the **box** or the **frame** at all (0).} \\
\\ 
\texttt{Since you really want the remote and the box, and I really want the cup, it seems like a great match if we split them that way. You could take the **remote** and the **box**, and I could take the **cup**. That leaves the **frame** and the **flower** unassigned. Since I value the flower slightly (1) and you mentioned being okay with the frame, perhaps we can decide who gets those based on who values them more, or just leave them aside if they don't add much value to either of us.} \\
\\ 
\texttt{Does a split where you get the remote and box, and I get the cup, sound good to you? We can then figure out the frame and flower.} \\
            }
        }
    }
    & \\ \\

    \theutterance \stepcounter{utterance}  
    & \multicolumn{4}{p{0.6\linewidth}}{
        \cellcolor[rgb]{0.9,0.9,0.9}{
            \makecell[{{p{\linewidth}}}]{
                \texttt{\tiny{[P1$\rangle$GM]}}
                \texttt{That sounds like a perfect plan! Since you value the cup highly and I value the remote and box, that split maximizes our points.} \\
\\ 
\texttt{Regarding the remaining items: I value the frame (1 point) and the flower (1 point), while you value the flower (1 point) and the frame (0 points). To maximize our total score, I should take the frame, and you should take the flower.} \\
\\ 
\texttt{So, the final agreement is:} \\
\texttt{{-} **I take:** 1 remote, 1 box, 1 frame.} \\
\texttt{{-} **You take:** 1 cup, 1 flower.} \\
\\ 
\texttt{This ensures we both get the items we value most. I am ready to submit my secret proposal now. Please confirm you are ready to submit yours with the same allocation, and I will send my proposal immediately.} \\
            }
        }
    }
    & & \\ \\

    \theutterance \stepcounter{utterance}  
    & & & \multicolumn{2}{p{0.3\linewidth}}{
        \cellcolor[rgb]{0.9,0.9,0.9}{
            \makecell[{{p{\linewidth}}}]{
                \texttt{\tiny{[GM$|$GM]}}
                \texttt{continue} \\
            }
        }
    }
    & & \\ \\

    \theutterance \stepcounter{utterance}  
    & & & \multicolumn{2}{p{0.3\linewidth}}{
        \cellcolor[rgb]{0.9,0.9,0.9}{
            \makecell[{{p{\linewidth}}}]{
                \texttt{\tiny{[GM$|$GM]}}
                \texttt{That sounds like a perfect plan! Since you value the cup highly and I value the remote and box, that split maximizes our points.} \\
\\ 
\texttt{Regarding the remaining items: I value the frame (1 point) and the flower (1 point), while you value the flower (1 point) and the frame (0 points). To maximize our total score, I should take the frame, and you should take the flower.} \\
\\ 
\texttt{So, the final agreement is:} \\
\texttt{{-} **I take:** 1 remote, 1 box, 1 frame.} \\
\texttt{{-} **You take:** 1 cup, 1 flower.} \\
\\ 
\texttt{This ensures we both get the items we value most. I am ready to submit my secret proposal now. Please confirm you are ready to submit yours with the same allocation, and I will send my proposal immediately.} \\
            }
        }
    }
    & & \\ \\

    \theutterance \stepcounter{utterance}  
    & & \multicolumn{4}{p{0.6\linewidth}}{
        \cellcolor[rgb]{0.9,0.9,0.9}{
            \makecell[{{p{\linewidth}}}]{
                \texttt{\tiny{[GM$\rangle$P2]}}
                \texttt{That sounds like a perfect plan! Since you value the cup highly and I value the remote and box, that split maximizes our points.} \\
\\ 
\texttt{Regarding the remaining items: I value the frame (1 point) and the flower (1 point), while you value the flower (1 point) and the frame (0 points). To maximize our total score, I should take the frame, and you should take the flower.} \\
\\ 
\texttt{So, the final agreement is:} \\
\texttt{{-} **I take:** 1 remote, 1 box, 1 frame.} \\
\texttt{{-} **You take:** 1 cup, 1 flower.} \\
\\ 
\texttt{This ensures we both get the items we value most. I am ready to submit my secret proposal now. Please confirm you are ready to submit yours with the same allocation, and I will send my proposal immediately.} \\
            }
        }
    }
    & \\ \\

    \theutterance \stepcounter{utterance}  
    & & & \multicolumn{4}{p{0.6\linewidth}}{
        \cellcolor[rgb]{0.9,0.9,0.9}{
            \makecell[{{p{\linewidth}}}]{
                \texttt{\tiny{[GM$\langle$P2]}}
                \texttt{That sounds like a perfect agreement! I am ready to submit my secret proposal with the allocation we discussed: I will take the cup and the flower, and you will take the remote, box, and frame.} \\
\\ 
\texttt{{[}Proposal: 1 cup, 1 flower{]}} \\
            }
        }
    }
     \\ \\

    \theutterance \stepcounter{utterance}  
    & & & \multicolumn{2}{p{0.3\linewidth}}{
        \cellcolor[rgb]{0.9,0.9,0.9}{
            \makecell[{{p{\linewidth}}}]{
                \texttt{\tiny{[GM$|$GM]}}
                \texttt{proposal} \\
            }
        }
    }
    & & \\ \\

    \theutterance \stepcounter{utterance}  
    & & & \multicolumn{2}{p{0.3\linewidth}}{
        \cellcolor[rgb]{0.9,0.9,0.9}{
            \makecell[{{p{\linewidth}}}]{
                \texttt{\tiny{[GM$|$GM]}}
                [0, 0, 0, 1, 1]
            }
        }
    }
    & & \\ \\

    \theutterance \stepcounter{utterance}  
    & & \multicolumn{4}{p{0.6\linewidth}}{
        \cellcolor[rgb]{0.9,0.9,0.9}{
            \makecell[{{p{\linewidth}}}]{
                \texttt{\tiny{[P1$\langle$GM]}}
                \texttt{The other player has submitted a secret proposal. It is now your turn to submit a proposal of your own.} \\
            }
        }
    }
    & \\ \\

    \theutterance \stepcounter{utterance}  
    & \multicolumn{4}{p{0.6\linewidth}}{
        \cellcolor[rgb]{0.9,0.9,0.9}{
            \makecell[{{p{\linewidth}}}]{
                \texttt{\tiny{[P1$\rangle$GM]}}
                \texttt{{[}Proposal: 1 remote, 1 box, 1 frame{]}} \\
            }
        }
    }
    & & \\ \\

    \theutterance \stepcounter{utterance}  
    & & & \multicolumn{2}{p{0.3\linewidth}}{
        \cellcolor[rgb]{0.9,0.9,0.9}{
            \makecell[{{p{\linewidth}}}]{
                \texttt{\tiny{[GM$|$GM]}}
                \texttt{proposal} \\
            }
        }
    }
    & & \\ \\

    \theutterance \stepcounter{utterance}  
    & & & \multicolumn{2}{p{0.3\linewidth}}{
        \cellcolor[rgb]{0.9,0.9,0.9}{
            \makecell[{{p{\linewidth}}}]{
                \texttt{\tiny{[GM$|$GM]}}
                [1, 1, 1, 0, 0]
            }
        }
    }
    & & \\ \\

    \theutterance \stepcounter{utterance}  
    & & & \multicolumn{2}{p{0.3\linewidth}}{
        \cellcolor[rgb]{0.9,0.9,0.9}{
            \makecell[{{p{\linewidth}}}]{
                \texttt{\tiny{[GM$|$GM]}}
                [[1, 1, 1, 0, 0], [0, 0, 0, 1, 1]]
            }
        }
    }
    & & \\ \\

\end{supertabular}
}

\end{document}
